\documentclass[10pt,a4paper]{amsart}
\usepackage[margin=19mm]{geometry}
\usepackage{amsmath,amssymb,mathtools}
\usepackage[scr=rsfs,cal=euler]{mathalfa}
\usepackage{xfrac}
\usepackage[unicode,hidelinks,bookmarks=false]{hyperref}
\newcommand{\ie}{i.e.\ }
\newcommand{\cf}{cf.\ }
\newcommand{\resp}{respectively\ }
\newcommand{\Subsection}{\subsection}
\providecommand{\nobreakdash}{\nobreak\hbox{-}}
\setlength{\emergencystretch}{2em}
\multlinegap0pt
\usepackage{amssymb}
\usepackage{xcolor}
\usepackage{mathtools}
\usepackage{xfrac}
\newtheorem{proposition}{Proposition}
\newtheorem{lemma}{Lemma}
\newtheorem{theorem}{Theorem}
\newcommand{\A}{\mathbb{A}}
\newcommand{\F}{\mathbb{F} }
\title{Generalized affine buildings for semisimple algebraic groups over real closed fields}
\author{Raphael Appenzeller}
\date{Source: arXiv:2601.03226v1 (CC BY 4.0)}
\begin{document}
\maketitle

\noindent\emph{Source and licence.} Generalized affine buildings for semisimple algebraic groups over real closed fields. Version arXiv:2601.03226v1 (CC BY 4.0), distributed by arXiv under the Creative Commons Attribution 4.0 International licence, http://creativecommons.org/licenses/by/4.0/, as recorded in the arXiv licence metadata for this exact version and shown on the abstract page https://arxiv.org/abs/2601.03226v1. Full licence notice: \texttt{licenses/arxiv-2601.03226-CC-BY-4.0.txt}.\par\smallskip

\noindent\textit{Editorial scope.} Counted unit: the article's proposition prop:A2\_b (source0.tex lines 2975--2977). The article's complete original proof of it is delivered with it (source0.tex lines 2978--3019, one \texttt{proof} environment of 42 lines). No mathematics is written, repaired or re-derived: the statement and the proof are the article's own lines, in the article's own notation, and no retained line is substituted or re-flowed.  Retained uncounted as the source-local context the proof needs: lines 1634--1639; lines 2936--2938; lines 2949--2954.  Nothing was omitted from any retained span, so the package carries no omission record.  No retained line contains a citation, so the wrapper carries no bibliography and no bibliography entry.  No retained line contains a figure environment, an image-inclusion command or drawing code, so no image is delivered and the package declares no asset dependency.  Editorial formatting only, in addition: an AMS \texttt{amsart} preamble that restates the article's own declarations and macros used by the retained lines, namely the environments \texttt{proposition}, \texttt{lemma}, \texttt{theorem} on separate counters, together with the standard packages those lines need.  The retained environments print in this document as Proposition 1, Lemma 1, Theorem 1 in order of appearance, whereas the article prints the same items with the numbering of its own full document.  The complete native source of the article as distributed in the arXiv e-print for this exact version is bundled unchanged as \texttt{sources/192213/source0.tex}.
\begin{proposition} \label{prop:UonA}
	For all $u\in U_\F$ and $a \in A_\F$
	$$
	ua.o \in \mathbb{A} \iff ua.o = a.o. 
	$$
\end{proposition}

		\begin{lemma}\label{lem:BTstab_dir_lemma}
			For $p,q\in \A$ there is an order on $\Sigma$ such that $U_q^+ \subseteq U_p^+$.
		\end{lemma}

		\begin{theorem}\label{thm:BTstab_fin}% (8.10)
			Let $\Omega \subseteq \A$ be a finite subset. Then the pointwise stabilizer of $\Omega$ satisfies
			$$
			\operatorname{Stab}_{G_\F}(\Omega) = \hat{P}_\Omega = \langle N_\Omega, U_{\alpha,\Omega} \colon \alpha \in \Sigma \rangle =  U_\Omega^+ U_\Omega^- N_\Omega.
			$$
		\end{theorem}

		\begin{proposition}\label{prop:A2_b}
			Let $g\in G_\F$ and $\Omega = g^{-1}\A \cap \A$. Then there exists an element $w\in W_a$ such that for all $p \in \Omega$, $g.p = w(p)$.
		\end{proposition}

		\begin{proof}
			We may assume that $\Omega \neq \emptyset$. For subsets $Y \subseteq \Omega$ we consider
			$$
			N_{g,Y} := \left\{ n \in \operatorname{Nor}_{G_\F}(A_\F) \colon  g.p = n.p \text{ for all } p\in Y \right\},
			$$
			so that our goal is to prove that $N_{g,\Omega} \neq \emptyset$.
			If $Y=\{p\}$ is a singleton, then $N_{g,\left\{p\right\}}$ is non-empty, since if $p=a.o$ and $g.p = b.o$ for $a,b\in A_\F$, then $ba^{-1} \in N_{\{p\}}$. We will now show by induction on the size of $Y$, that $N_{g,Y} \neq \emptyset$ for all \emph{finite} sets $Y \subseteq \Omega$.
			
			Let $Y\subseteq \Omega$ finite and $p \in \Omega$. Assume that there are $n_{Y} \in N_{g,Y}$ and $n_p \in N_{g,\{p\}} $, then $g^{-1}n_Y$ and $g^{-1}n_p $ lie in the pointwise stabilizers $\operatorname{Stab}_{G_\F}(Y)$ and $ \operatorname{Stab}_{G_\F}(\{p\})$. We may choose an order on $\Sigma$ such that $U_{Y}^+ \subseteq U_{\{p\}}^+$, by making sure that the defining chamber for the	order based at some point in $Y$ contains $p$, see Lemma \ref{lem:BTstab_dir_lemma}.
			Then, by Theorem \ref{thm:BTstab_fin},
			\begin{align*}
				n_Y^{-1} n_p &\in N_Y U_Y^- U_Y^+ U_{\{p\}}^+ U_{\{p\}}^- N_{\{p\}} 
				=  N_Y U_{Y}^- U_{\{p\}}^+ U_{\{p\}}^- N_{\{p\}}\\
				&=  N_Y U_{Y}^- U_{\{p\}}^- U_{\{p\}}^+ N_{\{p\}}
				\subseteq N_Y U^- U^+ N_{\{p\}}
			\end{align*}
			Let $n_Y' \in N_Y, n_p' \in N_{\{p\}}$ such that $(n_Y')^{-1}n_Y^{-1}n_pn_p' \in U^{-} U^{+} \cap \operatorname{Nor}_{G_\F}(A_\F)$. % actually $ = \{ \operatorname{Id}\}$.
			Since every element of $U^+$ stabilizes some affine chamber pointwise, so does every element of $U^-U^+ \cap \operatorname{Nor}_{G_\F}(A_\F)$, by Proposition \ref{prop:UonA}. The element of the affine Weyl group represented by $(n_Y')^{-1}n_Y^{-1}n_pn_p'$ thus acts trivially on some affine chamber, hence acts trivially on all of $\A$. 
			Therefore $[n_Yn_Y'] = [n_pn_p'] \in W_a$ and $ n_Y n_Y' \in N_{g,Y} \cap N_{g,\{p\}} = N_{g,Y\cup \{p\}}$, concluding the induction.
			
			Recall that the affine Weyl group $W_a = W_s \ltimes \A$ is isomorphic to the quotient $\operatorname{Nor}_{G_\F}(A_\F) / \operatorname{Cen}_{G_\F}(A_\F)$. Let $\pi \colon \operatorname{Nor}_{G_\F}(A_\F) \to W_s$ be the induced map to the finite group $W_s$. Let $Y_0 \subseteq \Omega$ be a finite subset, so that $|\pi(N_{g,Y_0})|$ is minimal. For any $p \in \Omega$,
			$$
			N_{g,Y_0 \cup \{p\}} = N_{g,Y_0} \cap N_{g,\{p\}} \subseteq N_{g,Y_0}
			$$
			and thus by minimality $\pi(N_{g,Y_0 \cup \{p\}}) = \pi(N_{g,Y_0})$. We claim that in fact $N_{g,Y_0 \cup \{p\}} = N_{g,Y_0}$ for every $p \in \Omega$.
			Pick $n_0\in N_{g,Y_0}$. 
			%We claim that $w:= [n_0] \in W_a$ is the element that satisfies $g.p=w(p)$ for all $p\in \A$. 
			We decompose $w := [n_0] =  (t_a,\pi(n_0)) \in W_a = \A \ltimes W_s$, so that $t_a$ is translation given by multiplication of some $a\in A_\F$. For any $p \in \Omega$, there exists $n' \in N_{g,Y_0 \cup \{p\}}$ such that $\pi(n_0) = \pi(n')$ and we decompose similarly $w' := [n'] = (t_{a'},\pi(n')) \in W_a$. Acting on some $q\in Y_0$, we have
			$$
			a.(\pi(n_0)(q)) = n_0.q = g.q = n'.q = a'.(\pi(n')(q)) = a'.(\pi(n_0)(q)),
			$$
			which implies $t_a = t_{a'}$ and thus $w= (t_a,\pi(n_0)) = (t_{a}', \pi(n')) = w'$. For any $p \in \Omega$ we thus have
			$$
			n_0.p = w(p) = w'(p) =  n'.p = g.p,
			$$
			so $n_0 \in N_{g, Y_0\cup\{p\}}$ as required. Finally
			$$
			N_{g,\Omega} %= \bigcap_{p \in \Omega} N_{g,\{p\}} 
			= \bigcap_{p \in \Omega} N_{g, Y_0\cup \{p\}} = N_{g,Y_0} \neq \emptyset,
			$$
			so taking any $n \in N_{g,Y_0}$ provides the required element $w = [n] \in W_a$ with $w(p) = g.p$ for all $p \in \Omega$.
		\end{proof}

\end{document}
